\documentclass[11pt,a4paper]{article}
\usepackage[T1]{fontenc}
\usepackage[utf8]{inputenc}
\usepackage{amsmath,amssymb,amsthm}
\usepackage[margin=1in]{geometry}
\usepackage[colorlinks=true,linkcolor=blue,urlcolor=blue]{hyperref}
\theoremstyle{plain}
\newtheorem{theorem}{Theorem}
\newtheorem{lemma}[theorem]{Lemma}
\newtheorem{proposition}[theorem]{Proposition}
\newtheorem{corollary}[theorem]{Corollary}
\theoremstyle{definition}
\newtheorem{remark}[theorem]{Remark}

\title{The empirical closed form for OEIS A224056, proved}
\author{Adrian Perez Fontelles\\ \small Independent researcher}
\date{1 September 2026}
\begin{document}
\maketitle

\begin{abstract}
OEIS A224056 carries an empirical closed form for $a(n)$ --- not a recurrence relating
consecutive terms but an explicit formula. It is true, and it is decidable rather than
empirical. The entry counts arrays subject to a condition on a bounded window of consecutive lines, so the lines of an array are the
vertices of a finite digraph, an array is a walk in it, and $a(n)$ is a walk count on
$S=30592$ vertices. Any function of the form $\sum_j p_j(n)\lambda_j^{\,n}$ satisfies the
monic linear recurrence whose characteristic polynomial is
$\prod_j(x-\lambda_j)^{\deg p_j+1}$; here that polynomial is $q(x)=(x-1)^{29}$, of degree
$29$. Two things are then checked exactly: that the walk count satisfies the same
recurrence from some index on, which the Cayley--Hamilton theorem reduces to a finite
computation, and that the two sequences agree at $29$ consecutive indices past that
point. A monic recurrence determines every later term from its predecessors, so agreement
there is agreement for good.
\end{abstract}

\noindent\small 2020 Mathematics Subject Classification. 05A15, 11B37, 15A18.\normalsize

\section{The conjecture}

OEIS A224056 is ``Number of nX7 0..2 arrays with rows and columns unimodal and antidiagonals nondecreasing.''. It has offset $1$ and begins
\[
239, 10511, 164781, 1442005, 9064541, 46853641, 212819499, 880290006,\ \dots
\]
The entry states, as a conjecture contributed by its author and never marked settled:
\begin{quote}
Empirical polynomial of degree 28 (see link above)

Empirical: a(n) = (1/1656632731900575744000000)*n\^{}28 + (1/23666181884293939200000)*n\^{}27 + (197/76740680486486016000000)*n\^{}26 + (191951/1880146671918907392000000)*n\^{}25 + (7946527/2256176006302688870400000)*n\^{}24 + (1503431/15041173375351259136000)*n\^{}23 + (6490149623/2517761630221841203200000)*n\^{}22 + (2869187099/49047304484841062400000)*n\^{}21 + (17244728801/14013515567097446400000)*n\^{}20 + (75135493547/3269820298989404160000)*n\^{}19 + (2155814659109/5162874156299059200000)*n\^{}18 + (5153254826101/860479026049843200000)*n\^{}17 + (207659181456407/2125889358476083200000)*n\^{}16 + (8712634636231/30369847978229760000)*n\^{}15 + (203016953797057/6197928158822400000)*n\^{}14 - (924979865456093/4601492117913600000)*n\^{}13 + (90702894734870589571/14197903929822412800000)*n\^{}12 - (1464845746168698967/43023951302492160000)*n\^{}11 + (75805769618145177203353/128749174272707788800000)*n\^{}10 - (21987491942680503592807/9196369590907699200000)*n\^{}9 + (26992062232273400774023/1094805903679488000000)*n\^{}8 - (24794932990162999879711/255454710858547200000)*n\^{}7 + (781558326336701679916651/1224053822863872000000)*n\^{}6 - (268038043963217961282893/102004485238656000000)*n\^{}5 + (30428503889127494898091/2380104655568640000)*n\^{}4 - (40855116344732385853/539706271104000)*n\^{}3 + (2318656934226364141/5936768982144)*n\^{}2 - (225397425874997/1115464350)*n - 3602498 for n>9
\end{quote}
As of the ``Last modified'' line on the live entry (Oct 03 2025 23:13:55, revision 10) this is still
recorded as empirical, and nothing on the entry records it as proved. Write
\[
f(n)\;=\;\frac{n^{28}}{1656632731900575744000000} + \frac{n^{27}}{23666181884293939200000} + \frac{197 n^{26}}{76740680486486016000000} + \frac{191951 n^{25}}{1880146671918907392000000} + \frac{7946527 n^{24}}{2256176006302688870400000} + \frac{1503431 n^{23}}{15041173375351259136000} + \frac{6490149623 n^{22}}{2517761630221841203200000} + \frac{2869187099 n^{21}}{49047304484841062400000} + \frac{17244728801 n^{20}}{14013515567097446400000} + \frac{75135493547 n^{19}}{3269820298989404160000} + \frac{2155814659109 n^{18}}{5162874156299059200000} + \frac{5153254826101 n^{17}}{860479026049843200000} + \frac{207659181456407 n^{16}}{2125889358476083200000} + \frac{8712634636231 n^{15}}{30369847978229760000} + \frac{203016953797057 n^{14}}{6197928158822400000} - \frac{924979865456093 n^{13}}{4601492117913600000} + \frac{90702894734870589571 n^{12}}{14197903929822412800000} - \frac{1464845746168698967 n^{11}}{43023951302492160000} + \frac{75805769618145177203353 n^{10}}{128749174272707788800000} - \frac{21987491942680503592807 n^{9}}{9196369590907699200000} + \frac{26992062232273400774023 n^{8}}{1094805903679488000000} - \frac{24794932990162999879711 n^{7}}{255454710858547200000} + \frac{781558326336701679916651 n^{6}}{1224053822863872000000} - \frac{268038043963217961282893 n^{5}}{102004485238656000000} + \frac{30428503889127494898091 n^{4}}{2380104655568640000} - \frac{40855116344732385853 n^{3}}{539706271104000} + \frac{2318656934226364141 n^{2}}{5936768982144} - \frac{225397425874997 n}{1115464350} - 3602498.
\]

\section{The count is a walk count}

The entry's defining condition is local: it is a condition on a bounded window of consecutive lines. Take as vertices the
admissible configurations of such a window and put an edge from $u$ to $v$ when $v$ may
follow $u$, which the condition decides by inspection. An admissible array is then exactly a
walk, so with $M$ the adjacency matrix of that digraph on $S=30592$ vertices, and $u,w$ the
vectors recording which windows may start and end an array,
\[
a(n)\;=\;u^{\!\top}M^{\,g(n)}w
\]
for the linear function $g$ that the entry's shape supplies. In particular $a$ is
$C$-finite: it satisfies some monic integer linear recurrence, though not necessarily a short
one.

\section{A closed form is a recurrence}

\begin{lemma}\label{lem:ann}
Let $f(m)=\sum_{j}p_j(m)\lambda_j^{\,m}$ with the $\lambda_j$ distinct and the $p_j$
polynomials, and put $q(x)=\prod_j(x-\lambda_j)^{\deg p_j+1}$. Then $q$ applied to the shift
operator annihilates $f$: writing $q(x)=x^{r}-\sum_{i=1}^{r}c_ix^{r-i}$,
\[
f(m)\;=\;\sum_{i=1}^{r}c_i\,f(m-i)\qquad\text{for every }m.
\]
\end{lemma}

\begin{proof}
The shift operator $E$ acts on $m\mapsto\lambda^{m}p(m)$ as
$(E-\lambda)\bigl(\lambda^{m}p(m)\bigr)=\lambda^{m+1}\bigl(p(m+1)-p(m)\bigr)$, and
$p(m+1)-p(m)$ has degree one less than $p$. So $(E-\lambda)^{\deg p+1}$ kills
$\lambda^{m}p(m)$, and the factors of $q$ commute.
\end{proof}

Here $q(x)=(x-1)^{29}$, of degree $29$; the closed form is a polynomial in $n$, so this is $(x-1)^{29}$, and the recurrence it encodes is
\begin{equation}\label{eq:rec}
a(n)\;=\;29\,a(n-1) - 406\,a(n-2) + 3654\,a(n-3) - 23751\,a(n-4) + 118755\,a(n-5) - 475020\,a(n-6) + 1560780\,a(n-7) - 4292145\,a(n-8) + 10015005\,a(n-9) - 20030010\,a(n-10) + 34597290\,a(n-11) - 51895935\,a(n-12) + 67863915\,a(n-13) - 77558760\,a(n-14) + 77558760\,a(n-15) - 67863915\,a(n-16) + 51895935\,a(n-17) - 34597290\,a(n-18) + 20030010\,a(n-19) - 10015005\,a(n-20) + 4292145\,a(n-21) - 1560780\,a(n-22) + 475020\,a(n-23) - 118755\,a(n-24) + 23751\,a(n-25) - 3654\,a(n-26) + 406\,a(n-27) - 29\,a(n-28) + a(n-29).
\end{equation}

\section{The proof}

\begin{lemma}\label{lem:crit}
With $q$ as above, $a$ satisfies \eqref{eq:rec} for all large $n$ if and only if
$u^{\!\top}M^{\,j}q(M)w=0$ for every $j\ge0$, and it is enough to check $j<S$.
\end{lemma}

\begin{proof}
Substituting the walk expression, the residual $a(n)-\sum_ic_ia(n-i)$ equals
$u^{\!\top}M^{\,g(n)-r}q(M)w$ up to the fixed linear reparametrisation $g$. For the bound,
$M^{S}$ is an integer combination of $I,M,\dots,M^{S-1}$ by Cayley--Hamilton, so the values
$u^{\!\top}M^{j}q(M)w$ satisfy the monic recurrence given by the characteristic polynomial of
$M$; $S$ consecutive zeros therefore force all later ones, and the last nonzero value pins the
threshold exactly.
\end{proof}

\begin{theorem}
$a(n)=f(n)$ for every $n\ge10$.
\end{theorem}

\begin{proof}
The vector $w'=q(M)w$ was formed by $29$ matrix--vector products in exact integer
arithmetic, and $u^{\!\top}M^{j}w'$ evaluated for $j=0,1,\dots$, again exactly. Those integers
vanish from the point corresponding to $n=39$ onwards, and the run was continued until the
working vector was identically zero, which settles every larger $j$ at once. So $a$ satisfies
\eqref{eq:rec} for every $n>38$. By Lemma~\ref{lem:ann} $f$ satisfies \eqref{eq:rec}
identically. Both were then evaluated in exact arithmetic at the $29$ consecutive indices
$n=39,\dots,67$ and agree at each. A monic recurrence of order $29$
determines $a(m)$ from $a(m-1),\dots,a(m-29)$, and for every $m\ge68$ both
$m>38$ and $m-29\ge39$ hold, so induction on $m$ gives $a(m)=f(m)$ for every
$m\ge39$.
Below $n=39$ the recurrence argument gives nothing, since \eqref{eq:rec} is only
established for $a$ from $n=39$ on. The remaining indices $n=10,\dots,38$ were
therefore checked one at a time, directly against the walk count in exact integer arithmetic,
and agree.
\end{proof}

The range is tight in the sense that was checked: at $n=9$ the closed form and the sequence differ, so no larger range is available.

\section{Verification}

Four checks, each independent of the algebra above.

First, the walk model reproduces every one of the $17$ terms the entry publishes,
exactly, in integer arithmetic. This is what ties the model to the entry: the model is built
by reading the entry's English, and a misreading gives a different digraph and different
counts.

Second, the closed form was evaluated in exact rational arithmetic --- not floating point ---
at every published term from $n=10$ on, and agrees with the entry's own DATA at each.

Third, the annihilation test of Lemma~\ref{lem:crit} was carried out exactly, and the model
and the closed form were then compared directly at sixty consecutive indices beyond
$n=10$, far past where the proof already settles the matter.

Fourth, the closed form was deliberately perturbed --- by a constant and by a term linear in
$n$ --- and each perturbed form was rejected by the same comparison, so the test is not one
that anything would pass.

\begin{thebibliography}{9}
\bibitem{oeis} The OEIS Foundation, \emph{The On-Line Encyclopedia of Integer
Sequences}, \url{https://oeis.org}, 2026. Sequence A224056.
\bibitem{stanley} R.~P.~Stanley, \emph{Enumerative Combinatorics}, Volume 1, 2nd ed.,
Cambridge University Press, 2011. (Transfer-matrix method, Section 4.7; $C$-finite sequences,
Section 4.1.)
\bibitem{kp} M.~Kauers and P.~Paule, \emph{The Concrete Tetrahedron}, Springer, 2011.
(Closed forms and linear recurrences with constant coefficients.)
\bibitem{horn} R.~A.~Horn and C.~R.~Johnson, \emph{Matrix Analysis}, 2nd ed., Cambridge
University Press, 2013.
\end{thebibliography}

\end{document}
